\chapter{Constructing Groups and Group Actions}

\section{Constructing Groups from Automorphisms}

\begin{proposition}[The construction $G:H$]
Let $G$ be a group and let
\[
    H\leq \Aut(G).
\]
For $h\in H$ and $g\in G$, write $g^h$ for the image of $g$ under the automorphism $h$.

Define
\[
    X=G:H=\{gh:g\in G,\ h\in H\}
\]
with multiplication
\[
    (g_1h_1)(g_2h_2)
    =
    \bigl(g_1g_2^{h_1^{-1}}\bigr)(h_1h_2).
\]
Then $X$ is a group. Moreover,
\[
    G\lhd X,
    \qquad
    X/G\cong H.
\]
\end{proposition}

\begin{proof}
Equivalently, as a set $X=G\times H$, and the multiplication is
\[
    (g_1,h_1)(g_2,h_2)
    =
    \bigl(g_1g_2^{h_1^{-1}},h_1h_2\bigr).
\]
The identity element is $(1,1)$, and
\[
    (g,h)^{-1}
    =
    \bigl((g^{-1})^h,h^{-1}\bigr).
\]
Associativity follows from the fact that $H$ acts on $G$ by automorphisms.

Identifying $G$ with $G\times 1$ and $H$ with $1\times H$, we get
\[
    X=GH,
    \qquad
    G\cap H=1.
\]
Also,
\[
    h^{-1}gh=g^h\in G
    \qquad(g\in G,\ h\in H),
\]
so $G\lhd X$. Finally, the quotient map $X\to X/G$ kills $G$ and identifies the quotient with $H$.
\end{proof}

\begin{remark}
The group $G:\Aut(G)$ is called the \textbf{holomorph} of $G$:
\[
    \operatorname{Hol}(G)=G:\Aut(G).
\]
\end{remark}

\section{Examples from Cyclic Groups}

\begin{example}[The group $\Z_7:\Z_2$]
Let
\[
    G=\Z_7=\gen{g}.
\]
Then
\[
    \Aut(G)\cong \Z_6.
\]

Let $b\in\Aut(G)$ have order $2$. Then $b$ acts by inversion:
\[
    g^b=g^{-1}.
\]
Hence
\[
    X=G:\gen{b}
    =
    \gen{g}:\gen{b}
\]
has presentation
\[
    X=
    \gen{g,b\mid |g|=7,\ |b|=2,\ g^b=g^{-1}}.
\]
Since
\[
    (gb)^2=gg^b=g g^{-1}=1,
\]
we get
\[
    X\cong \Z_7:\Z_2\cong D_{14}.
\]
\end{example}

\begin{example}[The group $\Z_7:\Z_3$]
Let
\[
    G=\Z_7=\gen{g}.
\]
Choose $z\in\Aut(G)$ of order $3$. For example, we may take
\[
    g^z=g^2,
\]
because $2$ has order $3$ in $(\Z/7\Z)^\times$.

Then
\[
    X=G:\gen{z}
    =
    \gen{g,z\mid |g|=7,\ |z|=3,\ g^z=g^2}.
\]
Thus
\[
    X\cong \Z_7:\Z_3.
\]
This group is nonabelian.
\end{example}

\begin{example}[General cyclic group of prime order]
Let
\[
    G=\gen{g}\cong \Z_p,
\]
where $p$ is prime. Then
\[
    \Aut(G)\cong \Z_{p-1}.
\]

Let $b\in\Aut(G)$ satisfy
\[
    |b|=\ell.
\]
Then
\[
    \ell\mid (p-1),
    \qquad
    \gen{b}\cong \Z_\ell.
\]
The corresponding semidirect product is
\[
    X=G:\gen{b}
    \cong
    \Z_p:\Z_\ell.
\]
More explicitly, if
\[
    g^b=g^\lambda,
\]
where $\lambda\in(\Z/p\Z)^\times$ has order $\ell$, then
\[
    X
    =
    \gen{
        g,b
        \mid
        |g|=p,\ |b|=\ell,\ g^b=g^\lambda
    }.
\]

If $\ell=2$, then $\lambda=-1$, and
\[
    X\cong D_{2p}.
\]

If $\ell=p-1$, then
\[
    X\cong \Z_p:\Z_{p-1}
    =
    \operatorname{AGL}(1,p),
\]
the affine general linear group of degree $1$ over $\F_p$.
\end{example}

\begin{remark}[The case $p=7$]
Since
\[
    \Aut(\Z_7)\cong \Z_6,
\]
the possible images of a cyclic complement $\Z_6$ in $\Aut(\Z_7)$ have order
\[
    1,\ 2,\ 3,\ 6.
\]
Thus one obtains the following groups:
\[
    \Z_7\times \Z_6,
    \qquad
    (\Z_7:\Z_3)\times \Z_2,
    \qquad
    (\Z_7:\Z_2)\times \Z_3,
    \qquad
    \Z_7:\Z_6.
\]
\end{remark}

\begin{remark}[General cyclic case]
In general, for
\[
    G=\Z_n,
\]
we have
\[
    \Aut(G)\cong \Unit(n).
\]
Thus every subgroup
\[
    H\leq \Unit(n)
\]
gives a group
\[
    X=\Z_n:H.
\]
\end{remark}

\section{The Case \texorpdfstring{$\Z_{35}$}{Z35}}

\begin{example}
Let
\[
    G=\Z_{35}\cong \Z_5\times \Z_7.
\]
Then
\[
    \Aut(G)
    \cong
    \Aut(\Z_5)\times \Aut(\Z_7)
    \cong
    \Z_4\times \Z_6.
\]
Write
\[
    \Aut(G)=\gen{b_1}\times\gen{b_2},
\]
where
\[
    |b_1|=4,
    \qquad
    |b_2|=6.
\]

For any subgroup $H\leq \Aut(G)$, we can form
\[
    X=\Z_{35}:H.
\]
Moreover,
\[
    X
    <
    (\Z_5:\Z_4)\times(\Z_7:\Z_6).
\]

Take
\[
    h=b_1^2b_2^3.
\]
Then $h$ has order $2$ and acts by inversion on both factors $\Z_5$ and $\Z_7$. Hence $h$ acts by inversion on $\Z_{35}$, and therefore
\[
    X_0=\Z_{35}:\gen{h}
    \cong
    D_{70}.
\]
\end{example}

\begin{remark}
If
\[
    G=D_{2n}\cong \Z_n:\Z_2,
\]
then
\[
    \Aut(G)
    \cong
    \Z_n:\Unit(n)
    =
    \Z_n:\Aut(\Z_n).
\]
\end{remark}

\section{Finite Field Example}

\begin{example}[Affine groups over finite fields]
Let
\[
    F=\operatorname{GF}(p^f)=\F_{p^f},
    \qquad
    q=p^f.
\]
Let
\[
    G=F^+.
\]
Then
\[
    G\cong \Z_p^f.
\]

Also,
\[
    F^\times\cong \Z_{q-1}.
\]
For each $h\in F^\times$, multiplication by $h$ induces an automorphism of the additive group $F^+$:
\[
    a\longmapsto a^h=ah.
\]
Indeed,
\[
    (a+b)^h
    =
    (a+b)h
    =
    ah+bh
    =
    a^h+b^h.
\]
Thus
\[
    F^\times\leq \Aut(F^+).
\]

Define
\[
    X=F^+:F^\times.
\]
Then
\[
    X\cong \Z_p^f:\Z_{p^f-1}.
\]
The subgroup $F^\times\leq \operatorname{GL}(f,p)$ is a \textbf{Singer cycle}.
\end{example}

\begin{proposition}[Automorphisms of the affine group]
Let
\[
    X=F^+:F^\times,
    \qquad
    F=\F_{p^f}.
\]
Then
\[
    \Aut(X)
    \cong
    X:\operatorname{Gal}(F/\F_p).
\]
Since
\[
    \operatorname{Gal}(F/\F_p)\cong \Z_f,
\]
we have
\[
    \Aut(X)
    \cong
    \bigl(\Z_p^f:\Z_{p^f-1}\bigr):\Z_f.
\]
Equivalently,
\[
    \Aut(X)
    \cong
    \bigl(\Z_p^f:\gen{g}\bigr):\gen{z},
\]
where
\[
    |g|=p^f-1,
    \qquad
    |z|=f,
    \qquad
    g^z=g^p.
\]
Here $z$ is the Frobenius automorphism
\[
    z:x\longmapsto x^p.
\]
\end{proposition}

\begin{proof}
    Put
    \[
        q=p^f,\qquad V=F^+,\qquad H=F^\times.
    \]
    Then
    \[
        X=V:H.
    \]
    We write elements of \(X\) as pairs
    \[
        (u,a),
        \qquad
        u\in F,\ a\in F^\times,
    \]
    with multiplication
    \[
        (u,a)(v,b)=(u+av,ab).
    \]
    Thus \(V\) is identified with \(\{(u,1):u\in F\}\), and \(H\) is identified with
    \[
        \{(0,a):a\in F^\times\}.
    \]
    
    We assume \(q>2\). The case \(q=2\) is exceptional, since then
    \[
        X\cong \Z_2,
        \qquad
        \Aut(X)=1.
    \]
    
    First observe that \(V\) is the unique Sylow \(p\)-subgroup of \(X\). Indeed,
    \[
        |V|=q=p^f,
        \qquad
        |H|=q-1,
    \]
    and \(p\nmid(q-1)\). Since \(V\lhd X\), it is the unique Sylow \(p\)-subgroup. Hence \(V\) is characteristic in \(X\).
    
    We first compute the normalizer of \(H=F^\times\) in
    \[
        \operatorname{GL}_{\F_p}(V).
    \]
    Here \(H\) acts on \(V=F^+\) by scalar multiplication:
    \[
        a:v\longmapsto av.
    \]
    We claim that
    \[
        N_{\operatorname{GL}_{\F_p}(V)}(H)
        =
        H:\operatorname{Gal}(F/\F_p).
    \]
    
    It is clear that \(H\) normalizes itself. Also, every
    \[
        \sigma\in \operatorname{Gal}(F/\F_p)
    \]
    normalizes \(H\), because
    \[
        \sigma(av)=\sigma(a)\sigma(v).
    \]
    Thus
    \[
        H:\operatorname{Gal}(F/\F_p)
        \leq
        N_{\operatorname{GL}_{\F_p}(V)}(H).
    \]
    
    Conversely, let
    \[
        A\in N_{\operatorname{GL}_{\F_p}(V)}(H).
    \]
    Let
    \[
        c=A(1)\in F^\times.
    \]
    Put
    \[
        B=c^{-1}A,
    \]
    where \(c^{-1}\) denotes scalar multiplication by \(c^{-1}\). Then \(B\) still normalizes \(H\), and
    \[
        B(1)=1.
    \]
    Since \(B\) normalizes \(H\), for each \(a\in F^\times\), there exists \(a^*\in F^\times\) such that
    \[
        B(ax)=a^*B(x)
        \qquad
        \text{for all }x\in F.
    \]
    Putting \(x=1\), we obtain
    \[
        B(a)=a^*.
    \]
    Therefore
    \[
        B(ax)=B(a)B(x)
        \qquad
        \text{for all }a\in F^\times,\ x\in F.
    \]
    This also holds for \(a=0\). Since \(B\) is \(\F_p\)-linear, it preserves addition:
    \[
        B(x+y)=B(x)+B(y).
    \]
    Thus \(B\) is a field automorphism of \(F\) fixing \(\F_p\). Hence
    \[
        B\in \operatorname{Gal}(F/\F_p).
    \]
    So
    \[
        A=cB\in H:\operatorname{Gal}(F/\F_p).
    \]
    This proves
    \[
        N_{\operatorname{GL}_{\F_p}(V)}(H)
        =
        H:\operatorname{Gal}(F/\F_p).
    \]
    
    Now let
    \[
        \alpha\in \Aut(X).
    \]
    Since \(V\) is characteristic in \(X\), we have
    \[
        \alpha(V)=V.
    \]
    The subgroup \(\alpha(H)\) is a complement to \(V\) in \(X\). We show that every complement to \(V\) is conjugate to \(H\) by an element of \(V\).
    
    Let \(K\) be a complement to \(V\). Then the projection
    \[
        X\to X/V\cong H
    \]
    restricts to an isomorphism
    \[
        K\cong H.
    \]
    Hence each element of \(K\) has the form
    \[
        (c(a),a),
        \qquad
        a\in H,
    \]
    for some map
    \[
        c:H\to V.
    \]
    Since \(K\) is a subgroup, we have
    \[
        (c(a),a)(c(b),b)=(c(ab),ab),
    \]
    and therefore
    \[
        c(ab)=c(a)+a c(b).
    \]
    Let \(r\) be a generator of \(H=F^\times\). Since \(q>2\), we have \(r\neq 1\). Define
    \[
        t=(1-r)^{-1}c(r)\in F.
    \]
    Then
    \[
        c(r)=t-rt.
    \]
    Using
    \[
        c(r^{n+1})=c(r^n)+r^n c(r),
    \]
    we get by induction
    \[
        c(r^n)=t-r^n t
        \qquad
        \text{for all }n.
    \]
    Thus
    \[
        K=(t,1)H(t,1)^{-1}.
    \]
    Therefore every complement to \(V\) is conjugate to \(H\) by an element of \(V\).
    
    Applying this to \(K=\alpha(H)\), there exists \(t\in V\) such that after composing \(\alpha\) with conjugation by \((t,1)^{-1}\), we may assume
    \[
        \alpha(H)=H.
    \]
    So it remains to describe the automorphisms of \(X\) that stabilize both \(V\) and \(H\).
    
    Assume now that
    \[
        \alpha(V)=V,
        \qquad
        \alpha(H)=H.
    \]
    Then
    \[
        \alpha|_V=A\in \operatorname{GL}_{\F_p}(V),
    \]
    and
    \[
        \alpha|_H=\theta\in \Aut(H).
    \]
    For \(a\in H\) and \(u\in V\), we have in \(X\)
    \[
        (0,a)(u,1)(0,a)^{-1}=(au,1).
    \]
    Applying \(\alpha\), we get
    \[
        (0,\theta(a))(A(u),1)(0,\theta(a))^{-1}
        =
        (A(au),1).
    \]
    Hence
    \[
        A(au)=\theta(a)A(u).
    \]
    Equivalently,
    \[
        A H A^{-1}=H.
    \]
    Thus
    \[
        A\in N_{\operatorname{GL}_{\F_p}(V)}(H).
    \]
    
    Conversely, if
    \[
        A\in N_{\operatorname{GL}_{\F_p}(V)}(H),
    \]
    then \(A\) induces an automorphism \(\theta_A\) of \(H\) by
    \[
        A\circ a\circ A^{-1}=\theta_A(a).
    \]
    Define
    \[
        \alpha_A:X\to X
    \]
    by
    \[
        \alpha_A(u,a)=(A(u),\theta_A(a)).
    \]
    Then
    \[
    \begin{aligned}
        \alpha_A\bigl((u,a)(v,b)\bigr)
        &=\alpha_A(u+av,ab)  \\
        &=(A(u+av),\theta_A(ab)) \\
        &=(A(u)+A(av),\theta_A(a)\theta_A(b)) \\
        &=(A(u)+\theta_A(a)A(v),\theta_A(a)\theta_A(b)) \\
        &=(A(u),\theta_A(a))(A(v),\theta_A(b)) \\
        &=\alpha_A(u,a)\alpha_A(v,b).
    \end{aligned}
    \]
    So \(\alpha_A\in\Aut(X)\).
    
    Therefore the subgroup of \(\Aut(X)\) stabilizing \(H\) is
    \[
        N_{\operatorname{GL}_{\F_p}(V)}(H)
        =
        H:\operatorname{Gal}(F/\F_p).
    \]
    
    The automorphisms induced by conjugation with elements of \(V\) form a subgroup isomorphic to \(V\). Indeed, if \((t,1)\in V\) centralizes \(H\), then
    \[
        t-at=0
        \qquad
        \text{for all }a\in F^\times.
    \]
    Since \(q>2\), we can choose \(a\neq 1\), and hence \(t=0\). Thus this copy of \(V\) intersects the stabilizer of \(H\) trivially.
    
    It follows that
    \[
        \Aut(X)
        \cong
        V:\bigl(H:\operatorname{Gal}(F/\F_p)\bigr).
    \]
    Since
    \[
        X=V:H=F^+:F^\times,
    \]
    we obtain
    \[
        \Aut(X)
        \cong
        X:\operatorname{Gal}(F/\F_p).
    \]
    
    Finally, for a finite field
    \[
        F=\F_{p^f},
    \]
    the Galois group is cyclic:
    \[
        \operatorname{Gal}(F/\F_p)
        =
        \langle z\rangle
        \cong
        \Z_f,
    \]
    where \(z\) is the Frobenius automorphism
    \[
        z:x\longmapsto x^p.
    \]
    Also,
    \[
        F^+\cong \Z_p^f,
        \qquad
        F^\times\cong \Z_{p^f-1}.
    \]
    Hence
    \[
        \Aut(X)
        \cong
        \bigl(\Z_p^f:\Z_{p^f-1}\bigr):\Z_f.
    \]
    
    If
    \[
        F^\times=\langle g\rangle,
    \]
    then
    \[
        |g|=p^f-1,
        \qquad
        |z|=f.
    \]
    The Frobenius automorphism acts on \(F^\times\) by
    \[
        a\longmapsto a^p.
    \]
    Therefore
    \[
        g^z=g^p.
    \]
    Thus equivalently
    \[
        \Aut(X)
        \cong
        \bigl(\Z_p^f:\langle g\rangle\bigr):\langle z\rangle,
    \]
    where
    \[
        |g|=p^f-1,
        \qquad
        |z|=f,
        \qquad
        g^z=g^p.
    \]
    \end{proof}

\begin{corollary}
For $p$ odd,
\[
    \Aut(\Z_p:\Z_{p-1})
    \cong
    \Z_p:\Z_{p-1}.
\]
\end{corollary}


    
    \begin{proof}
    Take
    \[
        F=\F_p.
    \]
    Then
    \[
        f=1,
        \qquad
        F^+\cong \Z_p,
        \qquad
        F^\times\cong \Z_{p-1}.
    \]
    Therefore
    \[
        X=F^+:F^\times
        \cong
        \Z_p:\Z_{p-1}.
    \]
    Since \(p\) is odd, we have \(p>2\), so the proposition applies.
    
    Moreover,
    \[
        \operatorname{Gal}(\F_p/\F_p)=1.
    \]
    Hence
    \[
        \Aut(X)
        \cong
        X:\operatorname{Gal}(\F_p/\F_p)
        \cong
        X.
    \]
    Therefore
    \[
        \Aut(\Z_p:\Z_{p-1})
        \cong
        \Z_p:\Z_{p-1}.
    \]
    \end{proof}

\section{The Case \texorpdfstring{$p=3$}{p=3} and \texorpdfstring{$f=2$}{f=2}}

\begin{example}
Assume
\[
    p=3,
    \qquad
    f=2.
\]
Then
\[
    q=9,
    \qquad
    F=\F_9.
\]
Let
\[
    X=\F_9^+:\F_9^\times.
\]
Since
\[
    \F_9^+\cong \Z_3^2,
    \qquad
    \F_9^\times\cong \Z_8,
\]
we have
\[
    X\cong \Z_3^2:\Z_8.
\]

By the previous proposition,
\[
    \Aut(X)
    \cong
    (\Z_3^2:\Z_8):\Z_2.
\]
Equivalently,
\[
    \Aut(X)
    \cong
    \Z_3^2:(\Z_8:\Z_2).
\]

Let
\[
    \Z_8=\gen{g},
\]
and let $z$ be the Frobenius automorphism. Since $p=3$, we have
\[
    g^z=g^3.
\]
Now define
\[
    a=g^2,
    \qquad
    b=gz.
\]
Then
\[
    |a|=4.
\]
Also,
\[
    b^2
    =
    (gz)(gz)
    =
    gzgz
    =
    gg^z z^2
    =
    gg^3
    =
    g^4.
\]
Hence
\[
    |b|=4,
    \qquad
    b^2=g^4.
\]
But
\[
    a^2=(g^2)^2=g^4,
\]
so
\[
    a^2=b^2.
\]

Moreover,
\[
    a^b
    =
    (g^2)^{gz}
    =
    (g^2)^z
    =
    g^6
    =
    (g^2)^{-1}
    =
    a^{-1}.
\]
Therefore
\[
    \gen{a,b}\cong Q_8,
\]
where
\[
    Q_8=\{1,-1,i,-i,j,-j,k,-k\}.
\]

Thus
\[
    \Aut(X)
    \cong
    (\Z_3^2:\Z_8):\Z_2
    =
    \Z_3^2:(\Z_8:\Z_2)
    >
    \Z_3^2:Q_8.
\]
\end{example}

\begin{exercise}
Show that all elements of order $3$ in
\[
    \Z_3^2:Q_8
\]
are conjugate.
\end{exercise}

\begin{proof}
    Let
    \[
        G=\Z_3^2:Q_8,
        \qquad
        V=\Z_3^2.
    \]
    Then
    \[
        V\lhd G,
        \qquad
        G/V\cong Q_8.
    \]
    
    First we show that every element of order \(3\) in \(G\) lies in \(V\). Let \(x\in G\) have order \(3\). Then the image \(xV\in G/V\cong Q_8\) has order dividing \(3\). But
    \[
        |Q_8|=8,
    \]
    so \(Q_8\) has no nontrivial element of order \(3\). Hence
    \[
        xV=V,
    \]
    and therefore
    \[
        x\in V.
    \]
    Since \(V\cong \Z_3^2\), every nonidentity element of \(V\) has order \(3\). Thus the elements of order \(3\) in \(G\) are exactly
    \[
        V^\#=V\setminus\{0\}.
    \]
    
    It remains to show that all elements of \(V^\#\) are conjugate in \(G\). Since \(Q_8\leq G\), it suffices to show that \(Q_8\) acts transitively on \(V^\#\).
    
    Recall that in this construction
    \[
        V=\F_9^+
    \]
    and
    \[
        Q_8=\langle a,b\rangle\leq \mathbb{F}_9^\times:\operatorname{Gal}(\mathbb{F}_9/\mathbb{F}_3),
    \]
    where
    \[
        a^2=b^2
    \]
    is the central element of order \(2\) in \(Q_8\). On \(V=\F_9^+\), this central element acts as multiplication by \(-1\). Hence it has no nonzero fixed point, because if
    \[
        -v=v,
    \]
    then
    \[
        2v=0,
    \]
    and over \(\F_3\) this implies
    \[
        v=0.
    \]
    
    Now let \(1\neq q\in Q_8\). If \(q\) has order \(2\), then \(q\) is the central element and we have just shown that it fixes no nonzero vector. If \(q\) has order \(4\), then
    \[
        q^2
    \]
    is the central element of order \(2\). If \(q\) fixed some \(0\neq v\in V\), then \(q^2\) would also fix \(v\), contradicting the previous paragraph. Therefore no nonidentity element of \(Q_8\) fixes a nonzero element of \(V\).
    
    Thus the stabilizer in \(Q_8\) of any \(0\neq v\in V\) is trivial. By the orbit-stabilizer theorem,
    \[
        |v^{Q_8}|
        =
        |Q_8:C_{Q_8}(v)|
        =
        |Q_8|
        =
        8.
    \]
    But
    \[
        |V^\#|=|\Z_3^2|-1=9-1=8.
    \]
    Therefore
    \[
        v^{Q_8}=V^\#.
    \]
    So \(Q_8\) is transitive on \(V^\#\).
    
    Consequently all nonidentity elements of \(V\), equivalently all elements of order \(3\) in \(G\), are conjugate in \(G\).
    \end{proof}

\section{Group Actions}

\begin{definition}[Group action]
Let $G$ be a group and let $\Omega$ be a set. We say that $G$ \textbf{acts} on $\Omega$ if the following conditions hold.

\begin{enumerate}[label=(\arabic*)]
    \item Each element $g\in G$ gives a one-to-one map
    \[
        \Omega\longrightarrow \Omega.
    \]
    The image of $\omega\in\Omega$ under $g$ is denoted by
    \[
        \omega^g.
    \]

    \item The identity element acts trivially:
    \[
        \omega^1=\omega
        \qquad
        \text{for all }\omega\in\Omega.
    \]

    \item The action is compatible with multiplication:
    \[
        (\omega^{g_1})^{g_2}
        =
        \omega^{g_1g_2}
        \qquad
        \text{for all }\omega\in\Omega,\ g_1,g_2\in G.
    \]
\end{enumerate}
\end{definition}

\begin{example}[Linear action]
The group $\operatorname{GL}(n,F)$ acts on
\[
    \Omega=F^n.
\]
Viewing vectors as row vectors, the action is
\[
    v^A=vA
    \qquad
    (v\in F^n,\ A\in\operatorname{GL}(n,F)).
\]
\end{example}

\begin{example}[Conjugation action]
A group $G$ acts on itself by conjugation:
\[
    \Omega=G,
    \qquad
    \omega^g=g^{-1}\omega g.
\]
The conjugacy class of an element $x\in G$ is
\[
    x^G
    =
    \{x^g:g\in G\}
    =
    \{g^{-1}xg:g\in G\}.
\]
\end{example}

\begin{example}[Right multiplication action]
A group $G$ acts on itself by right multiplication:
\[
    \Omega=G,
    \qquad
    \omega^g=\omega g.
\]
\end{example}

\section{The \texorpdfstring{$N/C$}{N/C} Lemma}

\begin{lemma}[$N/C$-lemma]
Let $H\leq G$. Then
\[
    C_G(H)\lhd N_G(H),
\]
and
\[
    N_G(H)/C_G(H)\leq \Aut(H).
\]
In particular, if $H\lhd G$, then
\[
    G/C_G(H)\leq \Aut(H).
\]
\end{lemma}

\begin{proof}
The normalizer $N_G(H)$ acts on $H$ by conjugation:
\[
    h^x=x^{-1}hx
    \qquad
    (h\in H,\ x\in N_G(H)).
\]
Since $x\in N_G(H)$, we have $x^{-1}Hx=H$, so this conjugation map is an automorphism of $H$.

Thus we get a homomorphism
\[
    \theta:N_G(H)\longrightarrow \Aut(H),
    \qquad
    x\longmapsto (h\mapsto h^x).
\]
Its kernel is
\[
    \Ker\theta
    =
    \{x\in N_G(H):h^x=h\text{ for all }h\in H\}
    =
    C_G(H).
\]
Therefore
\[
    C_G(H)\lhd N_G(H),
\]
and by the First Isomorphism Theorem,
\[
    N_G(H)/C_G(H)
    \cong
    \operatorname{Im}\theta
    \leq
    \Aut(H).
\]

If $H\lhd G$, then
\[
    N_G(H)=G,
\]
so
\[
    G/C_G(H)\leq \Aut(H).
\]
\end{proof}

